\section{主观责任阻却事由}
\begin{enumerate}
  \item 责任年龄。
        \begin{enumerate}
          \item 不满 12。
          \item 12 - 14。新增，没考过。\al{2019 大连 13 岁男孩强奸}。
                \begin{enumerate}
                  \item 2 个行为。杀人 \hz 伤害行为。
                  \item 2 个结果。致人死亡 \hz 重伤（特别残忍手段 + 严重残疾）。\al{泼硫酸}。
                  \item 1 个程序。最高检核准。221。
                \end{enumerate}
          \item \bs{14 - 16}。\zd{8 种故意犯罪行为}。
                故意\zd{杀}人、故意\zd{伤}害致人重伤或者死亡、强\zd{奸、抢}劫、\zd{投}放危险物质、贩卖\zd{毒}品、放\zd{火、爆}炸。
                \begin{enumerate}
                  \item 根据犯罪能力、作案能力筛选的。所以\zd{毒品仅限于贩卖毒品}。
                  \item \zd{帮助行为不负刑事责任}。
                \end{enumerate}
          \item 16。
        \end{enumerate}
  \item 责任能力。考察精神正不正常，鉴别能力。
        \begin{enumerate}
          \item 完全责任能力、相对/限定责任能力（半精神病）、完全无责任能力。
          \item 心理疾病 \bdy 精神病。程度不同，\zd{还有自制力}。
          \item 特殊人群。
                \begin{enumerate}
                  \item \al{间接性精神病患者}。
                  \item \al{盲} \yjt 必须从宽。
                  \item \al{聋哑} \yjt 居其一可以从宽，且才是必须从宽。
                  \item \al{生理性醉酒（喝高了）}\yjt 要。
                  \item \al{吸毒的人产生的幻觉，不分第几次吸} \yjt 要。
                \end{enumerate}
        \end{enumerate}
  \item \zd{行为与责任同时存在原则}。\zd{结合因果关系考}。
        \begin{enumerate}
          \item 行为与故意同时存在。\al{甲枪杀乙，心生怜悯，收枪不慎走火 \yjt 故意杀人预备，过失致人死亡，想象竞合择一重}。考点同结果的提前实现。
          \item 行为与责任年龄同时存在。包括作为和不作为。\al{甲 14 岁生日枪杀乙，乙重伤至第 2 日}。
          \item 行为与责任能力同时存在。不要求结果发生时，也要有责任能力。\zd{看正常行为和死亡结果有无因果关系}。
                \begin{enumerate}
                  \item 前行为（正常）\yjt 后行为（不正常）\yjt 死亡。\al{甲寄毒酒时正常，乙收到时甲精神不正常 \yjt 既遂}。
                  \item \al{甲第一刀正常，重伤，第二刀精神病崩溃，捅了第三刀 \yjt 二因一果}。
                  \item \al{甲第一刀正常但没砍中，第二刀精神病崩溃，捅了第三刀 \yjt 故意杀人未遂}。
                  \item 例外。\zd{原因自由行为}。第一次 \yjt 过失。后续再犯 \yjt 故意 \yjt 要（\zd{恶意利用}）。
                \end{enumerate}
        \end{enumerate}
  \item \bs{期待可能性}。法律不强人所难，已\zd{不具备谴责性}。\al{父亲包庇儿子；妻子包庇丈夫}。
        \begin{enumerate}
          \item vs 不作为。不作为指的是犯罪时，期待可能性更多强调犯罪后。
          \item \al{甲乙夫妻回娘家，甲和父母争吵，突发精神疾病，杀死父母，乙妻清洗血衣}。
          \item \al{\xhx{嫌疑人/被告人}作伪证、包庇、窝藏}。
        \end{enumerate}
  \item \zd{违法性认识可能性/法律认识错误}。刑法的禁止性、违法性。
        \begin{enumerate}
          \item 自然犯。杀人、抢劫、放火 \yjt 一般都具有违法性认识的可能性。
          \item 法定犯。以违法某个行政法为前提 \yjt \zd{只能听信有权机关}。只需认识到部分条件即可。
          \item 知道自己在干什么事（这一步认识错是事实错）\yjt 刑法对此事的评价态度（这一步认识错是法律错）。
          \item \zd{不知法者不免责}。例外。没有违法性的认识可能性。
        \end{enumerate}
\end{enumerate}
